% ============================================================================
% Financial Time-Series MLOps — Technical Report
% ============================================================================
\documentclass[11pt,a4paper]{article}

% ── Packages ─────────────────────────────────────────────────────────────────
\usepackage[margin=2.5cm]{geometry}
\usepackage[utf8]{inputenc}
\usepackage[T1]{fontenc}
\usepackage{lmodern}
\usepackage{amsmath}
\usepackage{booktabs}
\usepackage{longtable}
\usepackage{graphicx}
\usepackage{hyperref}
\usepackage{listings}
\usepackage{xcolor}
\usepackage{enumitem}
\usepackage{float}

% ── Listings style ───────────────────────────────────────────────────────────
\definecolor{codebg}{RGB}{245,245,245}
\definecolor{codeframe}{RGB}{200,200,200}
\definecolor{keyword}{RGB}{0,0,180}
\definecolor{comment}{RGB}{80,120,80}
\definecolor{string}{RGB}{180,0,0}

\lstset{
  basicstyle=\ttfamily\footnotesize,
  backgroundcolor=\color{codebg},
  frame=single,
  rulecolor=\color{codeframe},
  keywordstyle=\color{keyword}\bfseries,
  commentstyle=\color{comment}\itshape,
  stringstyle=\color{string},
  breaklines=true,
  breakatwhitespace=false,
  showstringspaces=false,
  tabsize=4,
  captionpos=b,
  numbers=left,
  numberstyle=\tiny\color{gray},
  numbersep=8pt,
}

\lstdefinestyle{python}{language=Python}
\lstdefinestyle{yaml}{language={},%
  morekeywords={stages,cmd,services,image,environment,ports,volumes,depends_on,build,healthcheck}
}
\lstdefinestyle{ascii}{language={},numbers=none}

% ── Metadata ─────────────────────────────────────────────────────────────────
\title{%
  \textbf{Financial Time-Series MLOps}\\[6pt]
  \large Predicting Next-Day Market-Open Direction\\
  for NSE F\&O Stocks Using OHLCV and News Sentiment
}
\author{
  Student Name\\
  Roll Number\\[4pt]
  \textit{MLOps Course End-Term Project}
}
\date{\today}

% ============================================================================
\begin{document}
\maketitle
\tableofcontents
\newpage

% ============================================================================
\section{Problem Statement}
% ============================================================================

This project builds an end-to-end MLOps pipeline to predict the
\textbf{next-day market-open price direction} (Bullish / Bearish / Neutral)
for 50 NSE Futures \& Options (F\&O) constituent stocks.

The prediction target is a three-class variable derived from the return
between the next day's 09:15~AM opening price and its closing price.
A return exceeding $+0.3\%$ is labelled \emph{Positive}, below $-0.3\%$ is
\emph{Negative}, and anything in between is \emph{Neutral}.

The system ingests two complementary data streams:
\begin{enumerate}[nosep]
  \item \textbf{OHLCV price data} — 15-minute intraday bars from TradingView,
        aggregated into daily candles with rolling technical indicators.
  \item \textbf{Financial news headlines} — scraped from Groww using Crawl4AI,
        enriched with FinBERT (ProsusAI/finbert) sentiment probabilities and
        768-dimensional CLS embeddings, then PCA-reduced and aggregated with
        exponential-decay rolling windows.
\end{enumerate}

The merged feature set (95 features) feeds an XGBoost / LightGBM classifier
trained with TimeSeriesSplit cross-validation and tracked via MLflow.
The best model is served through a containerised FastAPI endpoint and monitored
with Prometheus metrics.

\paragraph{Minimum expected scope covered.}
\begin{itemize}[nosep]
  \item Automated pipeline using \textbf{Airflow} (inference and training DAGs).
  \item Additional data-engineering tool: \textbf{PySpark} (headline cleaning, pre-open merge).
  \item Data preprocessing workflow (scraping $\to$ Spark cleaning $\to$ FinBERT $\to$ feature engineering).
  \item Training and comparison of \textbf{two models} (XGBoost and LightGBM).
\end{itemize}


% ============================================================================
\section{Dataset Description}
% ============================================================================

All data is stored as DVC-tracked CSV files in the repository's \texttt{data/}
directory.

\subsection{Stock Universe}

The universe comprises \textbf{50 NSE F\&O stocks} listed in
\texttt{data/stock\_list.csv}. Each row maps a company name
(\texttt{Stock\_name}) to its TradingView identifier, Groww slug, and NSE
symbol.

\subsection{OHLCV Price Data}

\begin{table}[H]
\centering
\caption{OHLCV data files}
\begin{tabular}{@{}llp{7cm}@{}}
\toprule
\textbf{File} & \textbf{Granularity} & \textbf{Description} \\
\midrule
\texttt{merged\_ohlc\_15min.csv} & 15-minute & Intraday bars for all 50 stocks.
  Columns: \texttt{symbol}, \texttt{datetime}, \texttt{open}, \texttt{high},
  \texttt{low}, \texttt{close}, \texttt{volume}. \\
\texttt{ohlc\_data/<stock>.csv} & 15-min / daily & Per-stock raw files (daily
  and 15-min), merged into the above. \\
\texttt{preopen\_csv/} & Point-in-time & NSE F\&O pre-open auction data,
  one dated CSV per trading day. \\
\bottomrule
\end{tabular}
\end{table}

OHLCV data is fetched via the TradingView data feed (\texttt{tvDatafeed}
library). Each stock has up to 10\,000 historical 15-minute bars
($\sim$2 years). The scraper supports incremental updates, fetching only
bars newer than the latest timestamp in the existing CSV.

\subsection{News Headlines}

\begin{table}[H]
\centering
\caption{Headline data files}
\begin{tabular}{@{}lp{9cm}@{}}
\toprule
\textbf{File} & \textbf{Description} \\
\midrule
\texttt{stock\_news/headlines.csv} & Raw scraped headlines (Crawl4AI from Groww).
  Columns: \texttt{headline\_id}, \texttt{symbol}, \texttt{published\_at},
  \texttt{source}, \texttt{headline}, \texttt{article\_url},
  \texttt{scraped\_at}. \\
\texttt{headlines.csv} & Cleaned and deduplicated headlines (PySpark output, DVC-tracked). \\
\texttt{headlines\_enriched.csv} & FinBERT-enriched: \texttt{positive\_prob},
  \texttt{negative\_prob}, \texttt{neutral\_prob}, \texttt{sentiment\_score},
  \texttt{sentiment\_label}. \\
\texttt{headlines\_enriched.embeddings.npy} & 768-d FinBERT CLS embeddings
  (NumPy binary). \\
\texttt{news\_features.csv} & PCA-reduced (50-d) + 7-day rolling aggregated
  sentiment features per (symbol, date). \\
\bottomrule
\end{tabular}
\end{table}

\subsection{Data Contracts}

All inter-role data exchange is governed by dataclass schemas defined in
\texttt{shared/schemas/data\_contract.py}:

\begin{lstlisting}[style=python, caption={Core data-contract schemas}]
@dataclass
class PriceRecord:
    symbol: str       # e.g. "Reliance Industries"
    date: str         # ISO-8601 date
    open: float
    high: float
    low: float
    close: float
    volume: int

@dataclass
class HeadlineRecord:
    headline_id: str
    symbol: str
    published_at: str # ISO-8601 datetime
    source: str
    headline: str
    author: Optional[str] = None
    body_snippet: Optional[str] = None
\end{lstlisting}


% ============================================================================
\section{System Architecture}
% ============================================================================

The project is organised into three roles, each with its own directory,
dependencies, and Airflow DAGs:

\begin{table}[H]
\centering
\caption{Role decomposition}
\begin{tabular}{@{}clp{7.5cm}@{}}
\toprule
\textbf{Role} & \textbf{Directory} & \textbf{Responsibility} \\
\midrule
1 & \texttt{role1\_data\_engineering/} & Scraping, PySpark cleaning, Airflow
  orchestration, DVC data versioning \\
2 & \texttt{role2\_ml\_modeling/}      & Feature engineering, model training
  (XGBoost / LightGBM), MLflow experiment tracking \\
3 & \texttt{role3\_mlops\_devops/}     & FastAPI serving, Docker containerisation,
  Prometheus monitoring, CI/CD \\
\bottomrule
\end{tabular}
\end{table}

Shared configuration lives in \texttt{shared/config.py} (environment-variable
driven, timezone-aware) and data contracts in
\texttt{shared/schemas/data\_contract.py}.

\subsection{High-Level Data Flow}

\begin{lstlisting}[style=ascii, caption={End-to-end data flow}]
  Scrapers (role1)              PySpark (role1)           ML Pipeline (role2)
  +-----------------+          +----------------+        +-----------------+
  | headline_scraper|--CSV---->| clean_headlines|--CSV-->| finbert_infer.  |
  | ohlc_scraper    |--CSV--+  | merge_preopen  |--CSV-->| sentiment_feat. |
  | nse_preopen     |--CSV--+  +----------------+        | feature_eng.    |
  +-----------------+     |         |                     | train.py        |
                          +--DVC--->+                     +--------+--------+
                                                                  |
                   Serving (role3)                                 |
                   +------------------+                   .pkl artifacts
                   | FastAPI + Docker |<-----DVC----------(DVC-tracked)
                   | batch_score.py   |
                   | Prometheus       |
                   +------------------+
\end{lstlisting}

\subsection{Orchestration}

Two Airflow DAGs drive the system:

\begin{enumerate}[nosep]
  \item \textbf{Inference DAG} (\texttt{financial\_ts\_inference\_pipeline}) —
        daily at 08:30~IST: scrape $\to$ Spark cleanup $\to$ Docker inference
        $\to$ DVC push.
  \item \textbf{Training DAG} (\texttt{financial\_ts\_training\_pipeline}) —
        manual or periodic: FinBERT $\to$ sentiment features $\to$ feature
        engineering $\to$ parallel XGBoost/LightGBM sweeps $\to$ DVC push models.
\end{enumerate}

\subsection{Data and Model Versioning}

DVC (Data Version Control) tracks all large artefacts:
\begin{itemize}[nosep]
  \item \texttt{data/headlines.csv.dvc}, \texttt{data/merged\_ohlc\_15min.csv.dvc}
  \item \texttt{models/*.pkl.dvc} (XGBoost, LightGBM, PCA, encoder, feature columns)
\end{itemize}
The DVC remote is Google Drive (service-account authentication).
Model versions are git-tagged (\texttt{model-v1}, \texttt{model-v2}).


% ============================================================================
\section{Data Pipeline}
\label{sec:data-pipeline}
% ============================================================================

\subsection{Scrapers}

Four modules ingest raw data daily. All support incremental updates.

\begin{table}[H]
\centering
\small
\begin{tabular}{@{}lllp{5.8cm}@{}}
\toprule
\textbf{Module} & \textbf{Source} & \textbf{Output} & \textbf{Notes} \\
\midrule
\texttt{headline\_scraper} & Groww (Crawl4AI) & \texttt{stock\_news/headlines.csv}
  & Async headless browser; MD5-based \texttt{headline\_id} for dedup \\
\texttt{ohlc\_scraper} & TradingView & \texttt{merged\_ohlc\_15min.csv}
  & Up to 10k 15-min bars/stock; per-stock CSVs merged \\
\texttt{nse\_preopen\_scraper} & NSE API & \texttt{preopen\_csv/*.csv}
  & F\&O pre-open auction; exponential-backoff retries \\
\texttt{finbert\_inference} & ProsusAI/finbert & enriched CSV + \texttt{.npy}
  & Batched sentiment probs + 768-d CLS embeddings;
    $s = (p_+ - p_-)(1 - p_0)$ \\
\bottomrule
\end{tabular}
\end{table}

\subsection{PySpark Processing}

Two PySpark jobs run in \texttt{local[*]} mode.
\textbf{clean\_headlines} applies a UDF (lowercase, strip HTML, collapse
whitespace), deduplicates by \texttt{headline\_id} via a window function,
and writes \texttt{data/headlines.csv}.
\textbf{merge\_preopen} maps NSE tickers to stock names via a Spark join on
\texttt{stock\_list.csv}, creates synthetic 09:15 OHLCV bars from the pre-open
data, deduplicates, and appends to \texttt{merged\_ohlc\_15min.csv}.
A third job, \textbf{trim\_for\_inference}, filters the full CSVs to the
minimal date window (45-day OHLCV, 10-day headlines) needed by the Docker
inference container.

\subsection{Airflow Inference DAG}

The DAG \texttt{financial\_ts\_inference\_pipeline} runs daily at 08:30~IST
with seven tasks:
scrape headlines + OHLCV (parallel) $\to$ scrape pre-open $\to$
Spark clean + merge (parallel) $\to$ Spark trim $\to$
Docker inference (\texttt{batch\_score.py}, trimmed data mounted read-only)
$\to$ DVC push.
Each task runs its module via \texttt{subprocess.run()} for import isolation.
Setting \texttt{SKIP\_SCRAPERS=1} disables live scraping for testing.


% ============================================================================
\section{Model Development}
\label{sec:model-dev}
% ============================================================================

% TODO: Fill in with role2 content.
% Suggested subsections:
%   5.1 Sentiment Feature Engineering (PCA 768->50, 7-day rolling, decay)
%   5.2 OHLCV Feature Engineering (daily aggregation, technical indicators,
%       9:15 features, target creation)
%   5.3 Model Selection (XGBoost vs LightGBM)
%   5.4 Training Procedure (TimeSeriesSplit CV, 3 param sets each)
%   5.5 MLflow Experiment Tracking

\textit{This section will detail the FinBERT $\to$ PCA $\to$ sentiment
aggregation pipeline, OHLCV technical-indicator feature engineering,
XGBoost and LightGBM training with TimeSeriesSplit cross-validation, and
MLflow experiment tracking. To be completed by the ML modelling team member.}


% ============================================================================
\section{Deployment Strategy}
\label{sec:deployment}
% ============================================================================

% TODO: Fill in with role3 content.
% Suggested subsections:
%   6.1 Docker Image (CPU-only torch, baked FinBERT, non-root)
%   6.2 Docker Compose Stack (API + Postgres + Prometheus)
%   6.3 FastAPI Endpoints (/health, POST /predict, GET /predictions)
%   6.4 Batch Scoring (batch_score.py, Airflow integration)

\textit{This section will describe the Docker containerisation strategy,
FastAPI serving layer, batch scoring integration with Airflow, and the
Postgres prediction store. To be completed by the MLOps team member.}


% ============================================================================
\section{Monitoring Strategy}
\label{sec:monitoring}
% ============================================================================

% TODO: Fill in with role3 content.
% Suggested subsections:
%   7.1 Prometheus Metrics (latency, request count, class distribution,
%       confidence)
%   7.2 Scrape Configuration
%   7.3 Example PromQL Queries

\textit{This section will describe the Prometheus-based monitoring setup,
the four custom metrics exposed at \texttt{/metrics}, and example PromQL
queries. To be completed by the MLOps team member.}


% ============================================================================
\section{CI/CD Implementation}
\label{sec:cicd}
% ============================================================================

% TODO: Fill in with role3 content.
% Suggested subsections:
%   8.1 GitHub Actions Workflow (test -> build Docker image)
%   8.2 Test Strategy
%   8.3 Image Build and Push

\textit{This section will describe the GitHub Actions CI/CD pipeline: test
job (pytest across all three roles), Docker image build on main branch
merges, and planned enhancements (container boot test, /health probe).
To be completed by the MLOps team member.}


% ============================================================================
\section{Results and Discussion}
\label{sec:results}
% ============================================================================

% TODO: Fill in with role2 content.
% Suggested subsections:
%   9.1 Cross-Validation Results (XGBoost F1=0.6631, LightGBM F1=0.6572)
%   9.2 Per-Class Performance
%   9.3 Feature Importance
%   9.4 Discussion

\textit{This section will present the cross-validation results for both
XGBoost (best F1 = 0.6631) and LightGBM (best F1 = 0.6572), per-class
metrics, feature importance analysis, and discussion. To be completed by
the ML modelling team member.}


% ============================================================================
\section{Challenges Faced}
\label{sec:challenges}
% ============================================================================

\begin{enumerate}


  % TODO: Add more challenges from role2/role3 as needed.
\end{enumerate}


% ============================================================================
\section{Future Improvements}
\label{sec:future}
% ============================================================================

% TODO: Fill in with team input.
% Suggested items:
%   - Drift detection against training-time feature baseline
%   - Minimal frontend over GET /predictions
%   - Alembic database migrations
%   - Prometheus Pushgateway for batch job metrics
%   - Additional data sources (options chain, FII/DII flows)

\textit{This section will outline planned enhancements: data-drift detection,
a prediction frontend, Alembic migrations, additional data sources, and
Prometheus Pushgateway integration for batch-job metrics. To be completed
collaboratively.}


% ============================================================================
\end{document}
